\documentclass[10pt,a4paper]{amsart}
\usepackage[margin=19mm]{geometry}
\usepackage{amsmath,amssymb,mathtools}
\usepackage[scr=rsfs,cal=euler]{mathalfa}
\usepackage{xfrac}
\usepackage[unicode,hidelinks,bookmarks=false]{hyperref}
\newcommand{\ie}{i.e.\ }
\newcommand{\cf}{cf.\ }
\newcommand{\resp}{respectively\ }
\newcommand{\Subsection}{\subsection}
\providecommand{\nobreakdash}{\nobreak\hbox{-}}
\setlength{\emergencystretch}{2em}
\multlinegap0pt
\usepackage{float}
\usepackage{amssymb}
\usepackage{braket}
\usepackage{mathtools}
\usepackage{tikz}
\usepackage{xcolor}
\newtheorem{lemma}{Lemma}
\newcommand{\Z}{\mathbb{Z}}
\newcommand{\e}{\operatorname{e}}
\title{On the exponential sum over squarefree integers}
\author{Nicolas Robles, Alexandru Zaharescu, Dirk Zeindler}
\date{Source: arXiv:2609.03961v1 (CC BY 4.0)}
\begin{document}
\maketitle

\noindent\emph{Source and licence.} On the exponential sum over squarefree integers. Version arXiv:2609.03961v1 (CC BY 4.0), distributed by arXiv under the Creative Commons Attribution 4.0 International licence, http://creativecommons.org/licenses/by/4.0/, as recorded in the arXiv licence metadata for this exact version and shown on the abstract page https://arxiv.org/abs/2609.03961v1. Full licence notice: \texttt{licenses/arxiv-2609.03961-CC-BY-4.0.txt}.\par\smallskip

\noindent\textit{Editorial scope.} Counted unit: the article's lemma lem:majorant (source0.tex lines 456--468). The article's complete original proof of it is delivered with it (source0.tex lines 470--501, one \texttt{proof} environment of 32 lines). No mathematics is written, repaired or re-derived: the statement and the proof are the article's own lines, in the article's own notation, and no retained line is substituted or re-flowed.  No further source-local context is needed: every cross-reference of every retained line resolves inside the two delivered spans.  Nothing was omitted from any retained span, so the package carries no omission record.  No retained line contains a citation, so the wrapper carries no bibliography and no bibliography entry.  No retained line contains a figure environment, an image-inclusion command or drawing code, so no image is delivered and the package declares no asset dependency.  Editorial formatting only, in addition: an AMS \texttt{amsart} preamble that restates the article's own declarations and macros used by the retained lines, namely the environments \texttt{lemma} on separate counters, together with the standard packages those lines need.  The retained environments print in this document as Lemma 1 in order of appearance, whereas the article prints the same items with the numbering of its own full document.  The complete native source of the article as distributed in the arXiv e-print for this exact version is bundled unchanged as \texttt{sources/209330/source0.tex}.
\begin{lemma}\label{lem:majorant}
Let $M\ge1$.  There is a trigonometric polynomial
\begin{equation}\label{eq:G-def}
  G_M(x)=\sum_{|h|\le H_M}c_h\e(hx),
  \qquad H_M\le\frac M2,
\end{equation}
with real coefficients satisfying $c_{-h}=c_h$ and
$0\le c_h\le12\log 2M$, such that $G_M(x)\ge0$ and
\begin{equation}\label{eq:G-majorizes}
  \min\left(M,\frac1{2\|x\|}\right)\le G_M(x)
\end{equation}
for all real $x$.
\end{lemma}

\begin{proof}
For an integer $H\ge1$, the Fej\'er kernel
\[
  F_H(x):=\frac1H\bigg|\sum_{0\le n<H}\e(nx)\bigg|^2
  =\sum_{|h|<H}\left(1-\frac{|h|}H\right)\e(hx)
  =\frac1H\left(\frac{\sin\pi Hx}{\sin\pi x}\right)^2
\]
is nonnegative with nonnegative coefficients at most $1$.  If
$0<\|x\|\le1/(4H)$, then $H\|x\|\le1/4$, so
$|\sin\pi Hx|=\sin(\pi H\|x\|)\ge2H\|x\|$, while $|\sin\pi x|\le\pi\|x\|$;
hence $F_H(x)\ge4H/\pi^2$, and the same holds at $x\in\Z$ since
$F_H(0)=H$.

If $M<8$, take $G_M=8$.  If $M\ge8$, put $J=\lceil\log_2M\rceil$ and
$H_j=2^{j-2}$ for $2\le j\le J$.  With $\delta=\|x\|$ we claim that
\begin{equation}\label{eq:dyadic-majorant}
  \min\left(M,\frac1{2\delta}\right)
  \le2+\sum_{2\le j\le J}2^{j-1}\mathbf1_{\{\delta\le2^{-j}\}}.
\end{equation}
If $\delta>1/4$, the left side is less than $2$.  If
$2^{-k-1}<\delta\le2^{-k}$ with $2\le k<J$, the right side is
$2+\sum_{j=2}^k2^{j-1}=2^k$ and the left side is less than $2^k$.  If
$\delta\le2^{-J}$, the right side is $2^J\ge M$.  Since
$1/(4H_j)=2^{-j}$, the polynomial
\[
  G_M(x):=2+\frac{\pi^2}2\sum_{2\le j\le J}F_{H_j}(x)
\]
majorizes the right side of \eqref{eq:dyadic-majorant}.  It is
nonnegative, its coefficients are nonnegative and symmetric, each is at most
$2+\frac{\pi^2}2(J-1)\le2+5\log_2M\le12\log2M$, and its degree is less
than $H_J=2^{J-2}<M/2$.
\end{proof}

\end{document}
